\section{实践落地：如何把推理时引导做成可复现实验}

\subsection{实验设计的四个关键旋钮}

从课程内容往工程落地推进时，最重要的不是再发明一个公式，而是把几个核心旋钮系统化管理起来。

\begin{enumerate}
    \item \textbf{采样器选择}：DDPM、SDE、ODE 各自对应不同的速度与稳定性权衡
    \item \textbf{引导强度调度}：固定权重最简单，但常常不如随时间步变化的 schedule 稳定
    \item \textbf{验证器质量}：reward / verifier 是否可信，直接决定引导方向是否正确
    \item \textbf{候选筛选策略}：单次采样、best-of-$N$、重排序（reranking）会显著改变最终效果
\end{enumerate}

\begin{importantbox}{很多实验失败并不是模型不行，而是 verifier 不够好}
一旦 verifier 本身与人类偏好不一致，或者只会奖励一些容易作弊的局部模式，扩散模型就会学会``投机取巧''。这与 RLHF 中 reward hacking 的问题完全同构。
\end{importantbox}

\subsection{推荐的实验日志模板}

为了让不同方法能被公平比较，课程内容很适合落成统一的实验记录模板：

\begin{table}[H]
\centering
\begin{tabular}{p{0.18\textwidth} p{0.22\textwidth} p{0.42\textwidth}}
\toprule
实验项 & 需要记录的字段 & 说明 \\
\midrule
基础模型 & 预训练模型、分辨率、训练域 & 影响 Tweedie 估计质量与可迁移性 \\
采样设置 & 步数、采样器、随机种子 & 决定速度与结果波动范围 \\
引导设置 & reward 类型、权重、schedule & 决定条件满足度与真实性平衡 \\
验证方式 & 自动指标、人工偏好、失败样例 & 防止只看平均分，忽略模式性崩坏 \\
计算开销 & 单样本耗时、显存、best-of-$N$ 次数 & 便于和训练时方法做成本比较 \\
\bottomrule
\end{tabular}
\caption{推理时引导实验的最小记录模板}
\end{table}

\begin{knowledgebox}{为什么 ``失败案例库'' 很重要}
扩散模型的引导方法往往在平均案例上看起来不错，但在某些条件组合或高分辨率场景下会系统性失真。为失败样例单独建库，比单纯汇报均值更有助于下一轮研究。
\end{knowledgebox}

\subsection{从单奖励到多奖励组合}

实际应用 rarely 只有一个目标。图像既要符合文本，又要满足几何一致性、风格约束、编辑指令、物理约束等。于是总奖励常写成：

$$r_{\text{total}}(x) = \lambda_1 r_{\text{text}}(x) + \lambda_2 r_{\text{consistency}}(x) + \lambda_3 r_{\text{physics}}(x) + \cdots$$

这带来两个新问题：

\begin{itemize}
    \item 不同奖励尺度不同，直接相加会让某一项支配优化
    \item 各奖励目标可能彼此冲突，例如更强的几何约束可能压缩图像细节与多样性
\end{itemize}

\begin{warningbox}{多奖励组合最常见的错误：权重凭直觉乱调}
如果没有做 reward normalization、梯度裁剪或时间步分配，不同约束会相互干扰，结果既不真实也不满足条件。一个更稳妥的做法是让不同奖励在不同时间段接管主导权，例如前期强调全局结构，后期强调细节和一致性。
\end{warningbox}

\subsection{研究前沿：从 hand-crafted guidance 到 learned verifier}

讲座的隐含趋势是：短期内，Tweedie 近似配合手工奖励函数仍然是性价比最高的路径；但长期来看，更可能的方向是训练更强的 verifier 或 value model：

\begin{itemize}
    \item \textbf{Learned reward model}：让 verifier 更接近人类偏好或下游任务成功率
    \item \textbf{Temporal value model}：不只评估最终 $x_0$，还评估中间状态 $x_t$ 的未来潜力
    \item \textbf{Adaptive compute}：把更多推理步数分配给困难样本，而不是平均对待所有输入
    \item \textbf{Hybrid pipeline}：训练时控制模型负责大方向，推理时引导负责末端精修
\end{itemize}

\subsection{评测协议：不要只看单一视觉指标}

课程虽然以方法推导为主，但从研究方法论上看，推理时引导尤其需要多维评测。一个更完整的 protocol 至少应包含：

\begin{enumerate}
    \item \textbf{条件满足度}：逆问题重建误差、约束满足率、任务成功率
    \item \textbf{真实性}：FID、CLIP score、人工主观评分或任务特定 perceptual metric
    \item \textbf{稳定性}：不同随机种子下结果波动是否可控
    \item \textbf{成本}：每次采样的 wall-clock time、显存、best-of-$N$ 带来的额外开销
\end{enumerate}

\begin{warningbox}{只汇报 ``最好的一张图'' 会严重夸大方法效果}
推理时引导很容易通过多次采样挑出漂亮案例，但如果不报告平均表现、失败率和计算成本，就无法公平比较不同方法。test-time scaling 的优势与代价必须同时披露。
\end{warningbox}

\subsection{为什么这条路线对中小团队尤其重要}

讲者强调的一个现实判断是：在大模型训练门槛持续升高的背景下，test-time 方法给了中小实验室另一条竞争路径。

\begin{itemize}
    \item 不需要从头训练大型生成模型
    \item 可直接复用公开权重与开源采样器
    \item 研究重点转移到 verifier 设计、奖励构造和 compute allocation
    \item 更适合快速试错，因为每次实验的反馈周期比重新训练短得多
\end{itemize}

\begin{knowledgebox}{扩散模型社区可能复制 LLM 社区的演化}
LLM 已经走出一条很清晰的路径：预训练模型商品化后，真正拉开差距的是后训练、评测和 test-time orchestration。扩散模型很可能也会沿着类似路线演进，未来竞争重点会越来越偏向 verifier、tooling 和 system design。
\end{knowledgebox}

\subsection{本章小结}

把推理时引导真正做成研究或产品系统，关键在于统一实验记录、设计稳定的 reward 组合、并把 verifier 质量视为第一等公民。很多看似算法层的问题，本质上是实验设计和评测设计的问题。

